\subsection{Weighted Frames and Equivariant Maps into Simplices}\label{sec:frame-simplex-duality}

In this section, we will explore the duality between (continuous) weighted frames and certain equivariant maps into simplices.
In particular, consider the setting of a $G$-space $X$ where $G$ is finite.
Then $\PP(G)$, the space of probability measures of $G$, is just the $(\card{G}-1)$-simplex, since we can label the standard basis vectors of $\R^{\card{G}}$ as $e_g$ for $g \in G$ and then write
\[
    \Delta^{\card{G}} = \set{\sum_{g \in G} x_g e_g \in \R^{\card{G}} : \sum_{g \in G} x_g = 1}
\]
which gives us the bijection
\[
    \PP(G) \ni \mu \mapsto \sum_{g \in G} \mu(g) e_g \in \Delta^{\card{G}}.
\]
Because of this, I will treat measures $\mu \in \PP(G)$ as vectors in $\R^{\card{G}}$.
Similarly, a function $f : G \to \R$ is equivalently a vector $f \in \R^{\card{G}}$.
This gives us that
\[
    \int_G f d \mu
    = \sum_{g \in G} \mu(g) f(g)
    = \mu \cdot f
\]
where $\cdot$ is the standard Euclidean dot product.

Going forward, we will assume that $G$ acts freely on $X$, which means that $\mu_x = \angles{\mu_x}_y$ for any weighted frame $\mu_\bullet$ and $x,y \in X$ since $G_y = \set{e}$.
In particular, for some convergent sequence $x_n \to x \in X$, this means $\angles{\mu_{x_n}}_x \wto \angles{\mu_x}_x$ if and only if $\mu_{x_n} \wto \mu_x$.
But this means that for all $f \in \R^{\card{G}}$, $\mu_{x_n} \cdot f \to \mu_x \cdot f$, and thus $\mu_{x_n} \to \mu_x$ under the usual Euclidean topology.
Hence, continuity of the weighted frame $\mu_\bullet$ is equivalent to continuity of the map $x \mapsto \mu_x$.
In summary, a continuous weighted frame for $G$ acting on $X$ is simply a continuous map from $X$ into $\Delta^{\card{G}}$ when $G$ is finite and acts freely on $X$.

We also want (weakly) equivariant weighted frames, which in this case are weighted frames such that $\mu_{hx} (g) = \mu_x (\inv{h} g)$ (equivariant and weakly equivariant frames are the same when $G$ acts freely).
Now let $\varphi : G \to \Sym(G)$ be the injective group homomorphism $\varphi(g) = \inv{g} \cdot -$ where $\Sym(G)$ is the symmetric group on $G$, so we want $\mu_{hx}(g) = \mu_x (\varphi(h) g)$.
But $\Sym(G)$ acts naturally on $\Delta^{\card{G}}$ via permutation, so an equivariant weighted frame is really just a $G$-equivariant map from $X$ to $\Delta^{\card{G}}$ where $G$ acts on $x \in \Delta^{\card{G}}$ by $g \cdot x := \varphi(g) x$.
In this setting, we have the following theorem.


\begin{theorem}\label{thm:equi-map-iff-sum}
    A map $f : X \to \Delta^{\card{G}-1}$ is $G$-equivariant if and only if there is some map $\alpha : X \to [0,1]$ such that $\sum_{g \in G} \alpha(gx) = 1$ for all $x \in X$ and
    \[
        f(x) = \sum_{g \in G} \alpha(\inv{g} x) e_g.
    \]
    \begin{proof}
        Note that the reverse direction is trivial, so we just prove the forward direction.
        Since $\card{G} < \infty$, we must be able to write $f(x) = \sum_{g \in G} f_g(x) e_g$ for continuous $f_g : X \to [0,1]$ with $\sum_{g \in G} f_g(x) = 1$ at all $x \in X$.
        Then the equivariance of $f$ gives us
        \[
            \sum_{g \in G} f_g(hx) e_g
            = f(hx)
            = h f(x)
            = h \sum_{g \in G} f_g(x) e_g
            = \sum_{g \in G} f_{\inv{h}g}(x) e_g.
        \]
        Hence, for each $g,h \in G$ and $x \in X$ we have $f_g(hx) = f_{\inv{h}g} (x)$, so we can take $\alpha := f_e$.
        Observe that $\alpha$ is certainly continuous, and we conclude by computing
        \[
            \sum_{g \in G} \alpha(\inv{g} x)
            = \sum_{g \in G} f_e (\inv{g} x)
            = \sum_{g \in G} f_g (x)
            = f(x).
        \]
    \end{proof}
\end{theorem}


We would generally like to construct equivariant frames that are computationally efficient, so we make the following definition.

\begin{definition}
    The \emph{cardinality} of a weighted frame $\mu_\bullet$ is
    \[
        \card{\mu_\bullet} := \sup_{x \in X} \min \set{n \in \N: \exists x_1,...,x_n \in X, \alpha_1,...,\alpha_n \in \R, \mu_x = \sum_{i \leq n} \alpha_i \delta_{x_i} \tand \sum_{i \leq n} \alpha_i = 1}.
    \]
    In other words, $\card{\mu_\bullet}$ is the minimum number of Dirac measures needed to represent $\mu_x$ for all $x \in X$.
    We take $\card{\mu_\bullet} = \infty$ if $\mu_x$ cannot be written as a finite sum of Dirac measures for some $x \in X$.
\end{definition}

Observe then that $\card{\mu_\bullet} = k$ means that the image of $x \mapsto \mu_x$ lives in the $(k-1)$-skeleton of $\Delta^{\card{G} - 1}$, and notably is not fully contained in the $(k-2)$-skeleton.
Hence, we have the following proposition.

\begin{proposition}
    Let $X$ be a free $G$-space for a finite group $G$.
    Then there exists a continuous equivariant weighted frame $\mu_\bullet : X \to \PP(G)$ with $\card{\mu_\bullet} = k$ if and only if there exists a $G$-equivariant map $f : X \to_G (\Delta^{\card{G}-1})^{(k-1)}$.
\end{proposition}
